\section{The theorem and its dependency map}\label{ch12:sec:theorem}
The theorem below says that the conditions of Section~\ref{ch12:sec:iteration} are enough, provided the space is complete: a contraction of a nonempty complete metric space has exactly one fixed point, and repeating the map from any starting point finds it. The theorem is also known as Banach's fixed-point theorem, after the Polish mathematician Stefan Banach (1892--1945), who published it in 1922.
\par\noindent\begin{minipage}{\linewidth}
\begin{theorem}[Contraction theorem]\label{thm:banach}
Let $(X,d)$ be a nonempty complete metric space, and let $T:X\to X$ be a contraction with factor $q$, where $0\le q<1$. Then the following hold.
\begin{enumerate}
\item $T$ has exactly one fixed point $x_*$.
\item For every starting point $x_0\in X$, the iterates $x_{n+1}=T(x_n)$ converge to $x_*$, and for every $n\in\NN$
\begin{equation}\label{eq:apriori}
d(x_n,x_*)\le\frac{q^n}{1-q}\,d(x_1,x_0).
\end{equation}
\item For every point $z\in X$, whether or not it is an iterate,
\begin{equation}\label{eq:residual}
d(z,x_*)\le\frac{d(z,Tz)}{1-q}.
\end{equation}
\end{enumerate}
\end{theorem}
\end{minipage}\par

In \eqref{eq:apriori}, the power $q^0$ means $1$, also when $q=0$. So for $n=0$ the bound reads $d(x_0,x_*)\le d(x_1,x_0)/(1-q)$.

\subsection*{Four separate promises}
The theorem makes four promises, and each of them needs its own argument.
\begin{itemize}
\item \emph{Existence}: some point of $X$ is left unchanged by $T$.
\item \emph{Uniqueness}: no two different points are left unchanged by $T$.
\item \emph{Convergence}: from every starting point, the iterates approach the fixed point.
\item \emph{Error bounds}: after finitely many steps, we can say how far we still are from the fixed point, even though we do not know where it is.
\end{itemize}
None of these follows automatically from the others. Section~\ref{ch04:sec:unique} showed that existence and uniqueness need separate proofs. Existence does not give convergence either: the map $T(x)=1-x$ of Section~\ref{ch12:sec:iteration} has the fixed point $1/2$, yet its iterates from $x_0=0$ jump between $0$ and $1$ forever.

\pauseq{The map $T(x)=1-x$ on $\RR$ has a fixed point, but its iterates from $x_0=0$ do not converge. Which hypothesis of Theorem~\ref{thm:banach} does it fail?}{The space $\RR$ is nonempty and complete, and $T$ maps $\RR$ to itself. The contraction inequality is what fails. For all real $x$ and $y$, $|T(x)-T(y)|=|(1-x)-(1-y)|=|y-x|=|x-y|$, so $T$ never shrinks a distance at all. A factor $q<1$ would require $|x-y|\le q|x-y|$, which is false whenever $x\ne y$.}

\subsection*{Two ways to measure the error}
The last two parts of the theorem answer the practical question ``how close are we?'' in two different ways.

The bound \eqref{eq:apriori} is called the \emph{a priori estimate}, from the Latin for ``from what comes before.'' It uses only the factor $q$ and the length $d(x_1,x_0)$ of the first step. So after a single step, before the rest of the iteration has been computed, it tells us how many steps will be enough for a desired accuracy.

The bound \eqref{eq:residual} is called the \emph{residual estimate}. The number $d(z,Tz)$ is the \emph{residual} of $z$: it measures how far $z$ is from satisfying the equation $Tz=z$, and it is zero exactly when $z$ is a fixed point. The point $z$ may come from anywhere: it may be an iterate, a guess, or the output of some other method. The residual estimate turns one evaluation of $T$ into a guaranteed upper bound on the error of $z$. A guaranteed bound of this kind, computed from known quantities, is called an \emph{error certificate}.

Both bounds divide by $1-q$, which is positive because $q<1$. This positive gap is what keeps the bounds finite. When $q$ is close to $1$, the number $1/(1-q)$ is large and the bounds are weak; at $q=1$ they would make no sense at all. Section~\ref{ch12:sec:estimates} shows both bounds at work.

\subsection*{The plan of the proof}
Here is the whole proof in outline, before the details. Each step names the earlier facts it uses.
\begin{enumerate}
\item \emph{The steps shrink geometrically.} The contraction inequality and induction (Chapter~\ref{ch:induction}) show that the step from $x_n$ to $x_{n+1}$ has length at most $q^n$ times the length of the first step.
\item \emph{The iterates form a Cauchy sequence.} Proposition~\ref{ch11:prop:geometric-cauchy}, which rests on the triangle inequality and the geometric sum, bounds the distance between any two late iterates.
\item \emph{A limit exists in $X$.} Completeness (Section~\ref{ch11:sec:complete}) turns the Cauchy sequence into a convergent one. Call its limit $x_*$.
\item \emph{The limit is a fixed point.} The triangle inequality and the contraction inequality show that $Tx_*$ and $x_*$ are closer than every positive tolerance, so they are equal.
\item \emph{There is only one fixed point.} The contraction inequality, applied to two fixed points, shows that they coincide. This is the comparison argument of Section~\ref{ch04:sec:unique} in a new setting.
\item \emph{The error bounds hold.} The bound from Step~2, the equation $Tx_*=x_*$ from Step~4 and two more uses of the triangle inequality give \eqref{eq:apriori} and \eqref{eq:residual}.
\end{enumerate}
A list like this, which records what each step of a proof depends on, is a \emph{dependency map} of the proof. It tells you what you need to know before you read the proof, and it shows where each hypothesis is used. Here the prerequisites are induction, the triangle inequality, geometric sums, Cauchy sequences and completeness. That is all: no calculus and no further theory is needed. Chapter~\ref{ch:bridge} uses the same kind of map to approach a proof whose ingredients are not all familiar yet.
